\subsection{Reference Safety}
\label{sec:ref}

The mutable-aliasing rule is critical in the presence of algebraic data types (|enum| type): a reference into a variant payload is only sound while the enum still occupies that variant.
If two mutable references to the same enum can simultaneously be live, one can switch the variant under the other.

\subsubsection{Example}
Consider the following code:

\begin{Move}
enum E { V1(Capability), V2(FakedCapability) }

fun aliased_swap(a: &mut E, b: &mut E, faked: &mut E) {
    // faked is V2(FakedCapability(victim))
    match (a) {
        E::V1(c) => {
            mem::swap(b, faked);  // *a is now V2(..)
            admin_action(c);      // c still &Capability
        },
        _ => {},
    }
}
\end{Move}

A call of the form !aliased_swap(&mut x, &mut x, ...)! should normally be rejected by the bytecode verifier's borrow-graph analysis because |a| and |b| would alias.

This was a real issue discovered via the bug bounty program after the introduction of algebraic data types.
There was a bug in the bytecode verifier's borrow analysis which did not correctly identify borrow graph edges |e.f1| and |e.f2| from different variants to refer to the same physical field at the same offset.
This allowed creating mutable references to the same field offset in the ADT, eventually enabling a call to !aliased_swap! with aliased mutable references.

\subsubsection{Performing the Checks}

Performing reference safety for borrow semantics in the style of Move (or Rust) at execution time is a novel problem which, to the best of our knowledge, has not been previously attempted.
The challenge is that during construction of references, mutable reference aliasing \emph{is} allowed.
Only once those references are used, aliasing must be forbidden.
As an example consider the following code -- obviously valid code which should be allowed:

\begin{Move}
struct S {a: u64, b: u64}
let s: S;
let r1 = &mut s.a;
let r2 = &mut s.b;
*r1 = 1; *r2 = 2;
\end{Move}

\noindent In the bytecode we generate for this, this is actually broken down into:

\begin{masm}
\BORROWLOC s \\
\BORROWFIELD a \\
\BORROWLOC s   \quad\quad \leftarrow \textit{temporary aliasing} \\
\BORROWFIELD b \\
\end{masm}

\noindent Thus while we have a mutable reference to $s.a$ on the stack, we create temporarily a new mutable reference to the whole of $s$, overlapping with the existing one.
Only after we consume this reference and create $s.b$, aliasing has been eliminated.
This property makes a naive approach to reference safety checking, for example via reference counting, infeasible.

Our implementation maintains a \emph{shadow state} that runs in parallel to the operand stack and the frame stack of the VM.
Every value carried by the running program is mirrored by a tag in the shadow stack: either ``non-reference'' or ``reference $\rho$'', where $\rho$ is a fresh identifier minted at the borrowing site.
For every non-reference value the shadow state can hold --- a local, a globally borrowed resource, or the value sitting behind a reference parameter passed in by the caller --- we build, lazily, an \emph{access path tree}: a tree whose edges are field (or variant) offsets, and whose nodes record the set of references currently pointing at that sub-location.
Vector elements collapse onto a single abstract edge so that vector-heavy programs do not blow up the tree.
A reference $\rho$ is then nothing more than a pointer into some access path tree, tagged with its mutability.

Aliasing is policed by \emph{poisoning}.
Whenever a destructive write happens through a mutable reference at some node $n$ --- this includes the implicit destructive write a $\MOVELOC$ performs on the local it empties --- the shadow state walks the tree around $n$ and marks every \emph{other} reference reachable from there as poisoned: all immutable references at $n$, its ancestors and its descendants, and all mutable references at strict descendants.
The mutable reference doing the write itself survives, matching the way the static borrow checker accounts for the borrow.
Any subsequent attempt to read, write, or further borrow through a poisoned reference aborts the transaction with an invariant-violation error.

At each function-call boundary the access path subtrees reached through the reference parameters are additionally locked --- exclusively for mutable parameters, shared for immutable --- which gives the callee the same single-writer / multiple-reader invariant the bytecode verifier statically assumes.
References returned by the callee are required to be derived from one of its reference parameters and are translated back to the corresponding access path in the caller's frame; native functions that mint or return references must declare an explicit reference model so the shadow state stays consistent across the call.

The bytecode fragment above runs unchanged: the second $\BORROWLOC\ s$ creates a new mutable reference at the root, but it is consumed by $\BORROWFIELD\ b$ before any destructive write fires, so no poisoning happens and execution proceeds.
The variant-switch attack from the start of this section, in contrast, is caught: the call to !mem::swap! is a destructive write at the variant root, which is an ancestor of the still-live payload reference |c|, and so |c| is poisoned the moment !swap! returns --- the subsequent !admin_action(c)! fails before it ever reads the forged token.



%%% Local Variables:
%%% mode: latex
%%% TeX-master: "main"
%%% End:
